% GPQA-Diamond: action-level + semantic-level loop mechanics analysis
% Supports Finding 1 — Agent loops without a clear target
% Integrates: canonical action analysis, ARM semantic mode analysis,
%             entropy trajectory analysis, keyword markers, context rotting

\subsection{Action-Level and Semantic Mechanics of Trajectory Inflation}
\label{appendix:macro-gpqa-loop}

The Finding~1 claim in \S\ref{appendix:macro-gpqa} states that OpenClaw's open
loop converts deliberation into trajectory inflation, where context exhaustion
rather than reasoning error becomes the proximate cause of incorrect answers.
We now decompose the mechanism at two levels: the canonical action level (what
operations the harness performs) and the ARM semantic level (what cognitive
strategy each reasoning segment represents).
We then diagnose \emph{why} the model enters the spiral, \emph{where} the token
budget goes inside it, and \emph{how} context rotting degrades output quality
as the trajectory lengthens.

% ===========================================================================
\paragraph{The loop is internal to a single assistant turn.}
% ===========================================================================

OpenClaw trajectories consist of exactly two message-level steps regardless of
outcome: one user message (system prompt plus task) and one assistant message.
Correct trajectories produce a compact assistant response (mean 2.4K tokens);
wrong trajectories produce a single continuous output stream averaging 28K
tokens, with 34\% reaching the 32K context limit before the model emits a valid
boxed answer.
The canonical action space $\{\texttt{plan},\texttt{reason},\texttt{tool\_use},
\texttt{verify},\texttt{recover},\texttt{answer}\}$ classifies the trajectory
as a whole.
On correct trajectories, OpenClaw's terminal action is \texttt{answer} in 99\%
of cases.
On regressions, only 6\% terminate in \texttt{answer}: the model ends in
\texttt{verify} (58\%) or \texttt{plan} (22\%), cycling without committing.

To analyze what happens \emph{inside} the single-turn output, we use the ARM
(Agentic Reasoning Mode) taxonomy, which classifies every segment of the
continuous output into one of seven cognitive modes:
principle-based deduction (PD), systematic elimination (SE), intuitive
commitment (IC), uncertainty navigation (UN), recovery-replanning (RR),
computational verification (CV), and undifferentiated reasoning (reason).

% ===========================================================================
\paragraph{Where the token budget goes: ARM mode allocation.}
% ===========================================================================

Table~\ref{tab:gpqa-arm-budget} reports the fraction of trajectory segments
allocated to each ARM mode.
This answers the question: when OpenClaw inflates from 2.4K to 28K tokens,
which cognitive activities consume the additional budget?

\begin{table}[t]
\centering
\small
\caption{%
  ARM cognitive-mode segment allocation by harness and outcome on GPQA-Diamond.
  Each cell is the mean fraction of segments classified to that mode.
  $\Delta$ = correct $-$ wrong difference in percentage points.
  Arm entropy = mode diversity (higher = more modes used).
}
\label{tab:gpqa-arm-budget}
\setlength{\tabcolsep}{3pt}
\begin{tabular}{llrrrrrrrc}
\toprule
Harness & Outcome & PD & reason & IC & UN & RR & SE & CV & Arm ent. \\
\midrule
\multirow{2}{*}{OpenClaw}
 & correct & 0.519 & 0.386 & 0.052 & 0.017 & 0.018 & 0.007 & 0.000 & 0.548 \\
 & wrong   & \textbf{0.608} & \textbf{0.291} & \textbf{0.026} & \textbf{0.043} & 0.010 & 0.022 & 0.000 & \textbf{0.395} \\
 & $\Delta$ & +8.9 & $-$9.5 & $-$2.6 & +2.6 & $-$0.8 & +1.5 & 0.0 & $-$0.15 \\
\midrule
\multirow{2}{*}{OpenCode}
 & correct & 0.353 & 0.587 & 0.026 & 0.009 & 0.019 & 0.006 & 0.000 & 0.593 \\
 & wrong   & \textbf{0.498} & \textbf{0.450} & 0.000 & \textbf{0.041} & 0.010 & 0.000 & 0.000 & \textbf{0.468} \\
 & $\Delta$ & +14.5 & $-$13.7 & $-$2.6 & +3.2 & $-$0.9 & $-$0.6 & 0.0 & $-$0.13 \\
\midrule
\multirow{2}{*}{ZeroClaw}
 & correct & 0.280 & 0.674 & 0.023 & 0.010 & 0.020 & 0.001 & 0.000 & 0.636 \\
 & wrong   & \textbf{0.406} & \textbf{0.504} & 0.011 & \textbf{0.032} & \textbf{0.044} & 0.003 & 0.000 & \textbf{0.812} \\
 & $\Delta$ & +12.6 & $-$17.0 & $-$1.2 & +2.2 & +2.4 & +0.2 & 0.0 & \textbf{+0.18} \\
\bottomrule
\end{tabular}
\end{table}

Three regularities emerge from Table~\ref{tab:gpqa-arm-budget}.
First, \textbf{PD is the dominant consumer of the inflated budget}.
Across all three harnesses, wrong trajectories allocate 8.9 to 14.5~pp more of
their segments to principle-based deduction than correct trajectories do, at the
direct expense of undifferentiated \texttt{reason} and IC.
The model is not ``over-thinking'' in a diverse or exploratory way; it is
specifically over-deriving, re-stating chemical principles or physical laws
without transitioning to elimination or commitment.

Second, \textbf{computational verification (CV) is zero across ALL harnesses
and outcomes}.
No harness successfully induces the Qwen3.5-27B model to perform
tool-mediated numerical verification as a distinct cognitive mode.
OpenCode achieves tool use at the canonical action level (55\% of correct
trajectories), but the ARM classifier assigns the model's reflection on the
tool output to \texttt{reason} rather than a dedicated verification mode.
OpenClaw never reaches CV either, despite its tool-intent markers growing 28-fold
on wrong trajectories.
The model \emph{wants} to use tools but does not translate intent into an
executable verification step.

Third, \textbf{arm entropy direction is harness-specific}.
OpenClaw and OpenCode wrong trajectories have \emph{lower} mode diversity than
correct ones ($-$0.15 and $-$0.13 respectively): the model collapses into fewer
modes.
ZeroClaw wrong trajectories have \emph{higher} mode diversity (+0.18): the model
oscillates between PD, reason, UN, and RR without converging.
The same ARM marker (high arm entropy) therefore signals productive exploration
under ZeroClaw but unproductive mode oscillation under OpenClaw.

% ===========================================================================
\paragraph{Why the model enters the spiral: the PD sink.}
% ===========================================================================

Table~\ref{tab:gpqa-arm-transitions-full} reports ARM transition probabilities
for OpenClaw, split by outcome.
Because OpenClaw outputs are single continuous messages, transitions represent
successive cognitive segments within one assistant turn.

\begin{table}[t]
\centering
\small
\caption{%
  ARM cognitive-mode transition probabilities within OpenClaw trajectories,
  by outcome.
  Each cell = probability of transitioning from row mode to column mode,
  averaged over all trajectories in the bucket.
  Bold = transition differs by $>0.10$ between correct and wrong.
}
\label{tab:gpqa-arm-transitions-full}
\setlength{\tabcolsep}{2.5pt}
\begin{tabular}{llrrrrrr}
\toprule
From $\to$ To & Outcome & PD & reason & UN & IC & RR & SE \\
\midrule
\multirow{2}{*}{PD}
  & correct & 0.709 & 0.248 & 0.004 & 0.017 & 0.012 & 0.002 \\
  & wrong   & \textbf{0.851} & \textbf{0.113} & 0.010 & 0.011 & 0.004 & 0.010 \\
\midrule
\multirow{2}{*}{reason}
  & correct & 0.521 & 0.301 & 0.065 & 0.069 & 0.016 & 0.008 \\
  & wrong   & \textbf{0.748} & \textbf{0.168} & \textbf{0.044} & 0.032 & 0.008 & 0.000 \\
\midrule
\multirow{2}{*}{UN}
  & correct & 0.200 & 0.467 & 0.200 & 0.067 & 0.000 & 0.000 \\
  & wrong   & 0.275 & 0.450 & 0.200 & 0.025 & 0.050 & 0.000 \\
\bottomrule
\end{tabular}
\end{table}

Table~\ref{tab:gpqa-arm-transitions-full} reveals a PD sink: once the model
enters principle-based deduction on a hard instance, the probability of staying
in PD on the next segment rises from 0.71 (correct) to 0.85 (wrong).
Simultaneously, the escape route narrows: PD$\to$\texttt{reason} drops from
0.25 to 0.11, and \texttt{reason}$\to$IC (the commitment path) collapses.
Even when the model pauses deduction for generic reflection
(\texttt{reason}$\to$PD at 0.75 on wrong trajectories vs.\ 0.52 on correct),
it immediately re-enters deduction.
The cognitive state is trapped: each deduction segment increases the probability
of another deduction segment, and the exit probability toward commitment decays
toward zero.

\emph{Over-reasoning vs.\ dedicated tool calling.}
The PD sink constitutes a specific failure of \emph{over-reasoning without
grounding}.
The model does not lack information: on OpenClaw regression tasks, DirectLLM
solves the same questions correctly at 11.3K tokens with entropy 0.31.
The model possesses the knowledge but cannot convert it into a commit decision
because the harness provides no exit signal.
Tool calling is not ``dedicated'' but \emph{aspirational}: tool-intent markers
explode to 53 per trajectory (vs.\ 1.9 on correct), yet the model never reaches
CV and only 7\% of wrong trajectories contain an actual tool invocation (vs.\
22\% on correct).
The model repeatedly expresses the desire to compute something
(\texttt{``Let me compute\ldots''}, \texttt{``I need to check\ldots''}) but
cycles back to PD without executing or without integrating the result.

% ===========================================================================
\paragraph{Context rotting: entropy collapse as the trajectory lengthens.}
% ===========================================================================

Table~\ref{tab:gpqa-context-rotting} reports three measures of how output
quality degrades with token position.
\emph{Head entropy} is the mean per-token log-probability entropy over the first
quartile of tokens; \emph{tail entropy} is the same over the last quartile.
\emph{Low-entropy token share} is the fraction of tokens whose entropy falls
below the 25th percentile of the full trajectory distribution.
\emph{Top-1 probability} is the mean probability the model assigns to its
selected token.

\begin{table}[t]
\centering
\small
\caption{%
  Context rotting signals: how output entropy and token probability evolve
  from the beginning (head) to the end (tail) of the trajectory.
  $\Delta$H = tail entropy $-$ head entropy.
  Low-ent share = fraction of tokens in the bottom entropy quartile.
  Top-1 prob = mean probability of the model's selected token.
}
\label{tab:gpqa-context-rotting}
\setlength{\tabcolsep}{3.5pt}
\begin{tabular}{llrrrrr}
\toprule
Harness & Outcome & Head ent. & Tail ent. & $\Delta$H & Low-ent share & Top-1 prob \\
\midrule
\multirow{3}{*}{OpenClaw}
 & correct   & 0.251 & 0.215 & $-$0.036 & 0.620 & 0.908 \\
 & regression & 0.241 & 0.073 & \textbf{$-$0.167} & \textbf{0.865} & \textbf{0.968} \\
 & both wrong & 0.254 & 0.147 & $-$0.107 & 0.739 & 0.954 \\
\midrule
\multirow{2}{*}{OpenCode}
 & correct   & 0.269 & 0.211 & $-$0.059 & 0.579 & 0.895 \\
 & regression & 0.236 & 0.165 & $-$0.071 & 0.663 & 0.912 \\
\midrule
\multirow{2}{*}{ZeroClaw}
 & correct   & 0.239 & 0.179 & $-$0.060 & 0.645 & 0.915 \\
 & regression & 0.244 & 0.237 & \textbf{$-$0.007} & \textbf{0.543} & \textbf{0.879} \\
 & both wrong & 0.244 & 0.265 & \textbf{+0.021} & 0.520 & 0.895 \\
\bottomrule
\end{tabular}
\end{table}

Table~\ref{tab:gpqa-context-rotting} reveals a harness-conditional context
rotting signature.
OpenClaw regressions begin at head entropy 0.241 (similar to correct trajectories
at 0.251) but collapse to tail entropy 0.073, a drop of 0.167.
By the final quartile, 86.5\% of tokens fall in the bottom entropy quartile
and the model assigns 0.97 probability to its selected tokens: the output has
become near-deterministic, cycling through high-confidence restatements rather
than exploring alternatives or advancing toward a decision.
The entropy collapse is not gradual but catastrophic: the model reaches a
critical context depth where diversity collapses and output becomes repetitive
filler.

ZeroClaw shows the opposite pattern.
Its regression trajectories maintain stable entropy from head to tail
($\Delta$H = $-$0.007), and its both-wrong trajectories show \emph{rising}
entropy (+0.021): the model explores genuinely under uncertainty without
collapsing into repetition.
Its low-entropy token share is 0.54 (vs.\ 0.87 for OpenClaw regression), and
top-1 probability is lower than correct trajectories (0.88 vs.\ 0.92), indicating
that the model remains in a deliberative, non-deterministic state throughout.
The fixed template prevents context rotting not by improving reasoning quality
but by capping the trajectory length before the rotting regime begins.

OpenCode regressions show a mild rotting signal ($\Delta$H = $-$0.071, low-ent
share 0.66), consistent with its intermediate architecture: tool-grounded
transitions provide partial exit ramps from PD, but on tasks where tools are
inapplicable the model still drifts toward lower-entropy output, just within a
shorter absolute token budget.

% ===========================================================================
\paragraph{Quantitative failure signatures: keyword markers.}
% ===========================================================================

Table~\ref{tab:gpqa-keyword-failure} reports per-trajectory counts of
diagnostic text patterns that distinguish spiraling from converging trajectories.

\begin{table}[t]
\centering
\small
\caption{%
  Quantitative keyword-marker analysis by harness and outcome.
  All values are per-trajectory means.
  Looping = restatements of question/goal.
  Self-corr = retractions or direction changes.
  Uncertainty = explicit hedging (``might be'', ``possibly'', ``not sure'').
  N-gram rate = fraction of 4-grams appearing more than once.
  Tool intent = expressions indicating desire to use a tool without actual
  invocation.
}
\label{tab:gpqa-keyword-failure}
\setlength{\tabcolsep}{3pt}
\begin{tabular}{llrrrrrr}
\toprule
Harness & Outcome & Looping & Self-corr & Uncert. & N-gram rep. & Tool intent & Tokens \\
\midrule
\multirow{3}{*}{OpenClaw}
 & correct   &    2 &    4 &   1 & 0.13 &   2 &  2{,}372 \\
 & regression &   83 &  148 &  35 & 0.63 &  53 & 27{,}420 \\
 & both wrong &   99 &  127 &  38 & 0.54 &  42 & 21{,}075 \\
\midrule
\multirow{3}{*}{OpenCode}
 & correct   &    3 &    6 &   1 & 0.44 &   8 &  1{,}036 \\
 & regression &   19 &   44 &   3 & 0.53 &  11 &  3{,}672 \\
 & both wrong &   11 &   29 &   6 & 0.55 &  12 &  2{,}996 \\
\midrule
\multirow{3}{*}{ZeroClaw}
 & correct   &    3 &    6 &   1 & 0.13 &   7 &  2{,}101 \\
 & regression &    5 &   14 &   3 & 0.12 &   7 &  1{,}913 \\
 & both wrong &    5 &   16 &   2 & 0.12 &   7 &  2{,}628 \\
\bottomrule
\end{tabular}
\end{table}

Table~\ref{tab:gpqa-keyword-failure} shows that OpenClaw's wrong trajectories
are not merely long; they are \emph{qualitatively different} in lexical
structure.
Self-correction markers explode from 4 (correct) to 148 (regression): the model
retracts or changes direction roughly once every 185 tokens on average.
Yet these self-corrections do not converge to a correct answer: they represent
oscillation between incompatible reasoning chains rather than progressive
refinement.
The repeated n-gram rate nearly quintuples (0.13 to 0.63), confirming that the
output is lexically self-similar rather than informationally progressive.
Uncertainty markers rise 35-fold (1 to 35), but the uncertainty is never
resolved: the model expresses doubt, re-enters PD to resolve it, and emerges
still uncertain, creating an uncertainty$\to$PD$\to$uncertainty cycle.

Critically, ZeroClaw shows none of these signatures despite sharing the same
model.
Its keyword markers are nearly identical across outcomes: looping markers stay
at 3--5, self-correction at 6--16, and repeated n-gram rate at 0.12--0.13
regardless of correctness.
The fixed template prevents the lexical degradation that characterizes
OpenClaw's spirals, even though the ARM mode rates (Table~\ref{tab:gpqa-arm-budget})
confirm that ZeroClaw's wrong trajectories also over-allocate to PD.

% ===========================================================================
\paragraph{Failure mode taxonomy.}
% ===========================================================================

Table~\ref{tab:gpqa-failure-taxonomy} classifies every incorrect trajectory
into a mechanistic cause bucket based on the relationship between the harness
output and the DirectLLM output on the same task.

\begin{table}[t]
\centering
\small
\caption{%
  Failure cause taxonomy by harness.
  Cause buckets are mutually exclusive and assigned by comparing harness output
  against DirectLLM output on the same task.
}
\label{tab:gpqa-failure-taxonomy}
\setlength{\tabcolsep}{3pt}
\begin{tabular}{lrrr}
\toprule
Cause bucket & $n$ & Mean tokens & Mean entropy \\
\midrule
\multicolumn{4}{c}{\textit{OpenClaw} (67 failures)} \\
\quad answer\_format\_or\_extraction (context exhaustion) & 32 & 30{,}451 & 0.051 \\
\quad both\_wrong (genuine reasoning error)               & 31 & 21{,}075 & 0.168 \\
\quad other (tool, verification, overrun)                &  4 &  3{,}169 & 0.304 \\
\midrule
\multicolumn{4}{c}{\textit{OpenCode} (53 failures)} \\
\quad both\_wrong (genuine reasoning error)               & 29 &  2{,}996 & 0.317 \\
\quad answer\_format\_or\_extraction (context exhaustion) & 15 &  4{,}718 & 0.107 \\
\quad premature\_answer                                   &  5 &  1{,}467 & 0.437 \\
\quad other (tool, verification)                         &  4 &  2{,}783 & 0.362 \\
\midrule
\multicolumn{4}{c}{\textit{ZeroClaw} (37 failures)} \\
\quad both\_wrong (genuine reasoning error)               & 22 &  2{,}628 & 0.372 \\
\quad premature\_answer                                   &  5 &  2{,}242 & 0.373 \\
\quad answer\_changed\_valid                              &  3 &  1{,}990 & 0.369 \\
\quad other (extraction, churn, verification)            &  7 &  2{,}105 & 0.348 \\
\bottomrule
\end{tabular}
\end{table}

Table~\ref{tab:gpqa-failure-taxonomy} reveals three mechanistically distinct
failure regimes.

\textbf{Context exhaustion failure} (OpenClaw dominant, 32/67 failures).
The model enters a PD sink, self-correction markers explode, entropy collapses
to 0.05, and the trajectory hits the 32K context limit before a valid boxed
answer is emitted.
These are tasks DirectLLM answers correctly (regressions): the base model knows
the answer, but the harness loop destroys it.
The 30.5K mean tokens represent almost pure waste: the model is not exploring
or computing, it is cycling through low-entropy restatements.

\textbf{Genuine reasoning error} (all harnesses, dominant in ZeroClaw and
OpenCode).
The model reaches a wrong answer through normal-length deliberation (2--3K
tokens) with healthy entropy (0.17--0.37).
DirectLLM also fails these tasks: the base model lacks the knowledge or
reasoning capability.
No harness architecture can rescue these failures, and the trajectories show
no pathological signatures.
The entropy is \emph{higher} than on correct trajectories for both ZeroClaw
(0.37 vs.\ 0.25) and OpenCode (0.32 vs.\ 0.29), consistent with genuine
uncertainty during deliberation.

\textbf{Premature commitment} (OpenCode and ZeroClaw).
The model emits an answer after minimal deliberation (1.5--2.2K tokens) with
high entropy (0.37--0.44): it commits under uncertainty rather than continuing
to reason.
This is the inverse failure mode to context exhaustion: insufficient rather than
excessive deliberation.

% ===========================================================================
\paragraph{The loop is task-structure-dependent.}
% ===========================================================================

Table~\ref{tab:gpqa-loop-subdomain} confirms that looping is not uniform across
the benchmark.
Subdomains where the task affords an executable, computational path (High-energy
physics, Quantum Mechanics) show minimal looping and near-perfect accuracy.
Subdomains requiring irreducible verbal reasoning about molecular structure
(Organic Chemistry, Molecular Biology) concentrate the spiral trajectories.

\begin{table}[t]
\centering
\small
\caption{%
  OpenClaw looping markers and token inflation by GPQA-Diamond subdomain,
  sorted by mean looping markers descending.
}
\label{tab:gpqa-loop-subdomain}
\setlength{\tabcolsep}{4pt}
\begin{tabular}{lrrrr}
\toprule
Subdomain & $n$ & Loop markers & Mean tokens & Acc\% \\
\midrule
Organic Chemistry           & 72 &  59 & 17{,}564 &  49 \\
Astrophysics                & 13 &  50 &  7{,}900 &  77 \\
Molecular Biology           & 15 &  36 &  7{,}533 &  67 \\
Relativistic Mechanics      &  7 &  27 &  7{,}528 &  71 \\
Chemistry (general)         & 20 &  21 &  6{,}746 &  75 \\
Physics (general)           & 19 &  11 &  7{,}724 &  79 \\
Genetics                    &  4 &   4 &  2{,}710 &  50 \\
High-energy particle phys.  & 14 &   3 &  1{,}391 & 100 \\
Quantum Mechanics           & 25 &   2 &  2{,}041 &  68 \\
Electromagnetism \& Photonics & 6 &   0 &  6{,}236 &  83 \\
\bottomrule
\end{tabular}
\end{table}

The subdomain gradient in Table~\ref{tab:gpqa-loop-subdomain} supports the
PD-sink mechanism.
Organic Chemistry questions (stereochemistry, reaction mechanisms, aromatic
substitution) require multi-step verbal reasoning about molecular structure.
The model defaults to PD because there is no numerical path to compute; the
open loop permits unbounded PD chaining; entropy collapses; and context
exhaustion prevents answer emission.
Quantum Mechanics questions involve sequential algebraic structure that OpenClaw
can navigate with brief verification (2.0K tokens, 68\% accuracy), while
ZeroClaw's fixed plan-reason-answer template maps directly onto the
decomposition structure, yielding 100\% accuracy on the same subdomain.

% ===========================================================================
\paragraph{Why ZeroClaw and OpenCode avoid the spiral.}
% ===========================================================================

The same Qwen3.5-27B model exhibits the same cognitive tendency toward PD
over-expression on wrong trajectories under all three harnesses
(Table~\ref{tab:gpqa-arm-budget}).
Yet only OpenClaw produces the catastrophic entropy-collapse-plus-context-exhaustion
failure mode.
The difference is architectural.

\textbf{ZeroClaw} imposes a fixed template: plan, then reason for a fixed number
of steps, then answer.
Its action count distribution is near-constant (correct: mean 18.0, std 0.19;
wrong: mean 16.9, std 2.7), confirming that the template overrides any tendency
to re-enter planning.
The template caps the trajectory at $\sim$2.3K tokens regardless of outcome,
which is below the threshold where context rotting begins.
ZeroClaw wrong trajectories show elevated PD (0.41 vs.\ 0.28) and arm entropy
(0.81 vs.\ 0.64), but the entropy is high rather than collapsed: the model
oscillates between modes under uncertainty but is forced to commit before
degradation sets in.
The failure is a genuine reasoning error under a budget constraint, not a
context-exhaustion artifact.

\textbf{OpenCode} provides a natural exit ramp from PD through tool execution.
When the task affords a computable path (numerical or symbolic structure),
executing code forces a mode transition out of PD and the result anchors the
answer.
When the tool path is unavailable (Organic Chemistry), OpenCode shows the same
PD over-expression as OpenClaw (PD rate 0.53 on regressions vs.\ 0.34 on
correct), but its wrong trajectories terminate at 3.7K tokens because the
harness does not permit unbounded re-entry.
Its entropy drop is mild ($\Delta$H = $-$0.07, low-ent share 0.66) compared to
OpenClaw ($\Delta$H = $-$0.17, low-ent share 0.87).

% ===========================================================================
\paragraph{Mechanism summary: three necessary conditions for the spiral.}
% ===========================================================================

The trajectory-inflation spiral operates through three interacting conditions,
each necessary and jointly sufficient.

\begin{enumerate}[leftmargin=*, itemsep=2pt]
\item \textbf{Task-structural tool-path conflict.}
  The model encounters a task where tool use is structurally unavailable (e.g.,
  reasoning about chemical reaction mechanisms) but its system prompt instructs
  it to use tools actively.
  The model repeatedly expresses tool intent (53 markers per wrong trajectory)
  without producing useful execution.
  On tool-affording subdomains (Quantum Mechanics, High-energy physics), the
  same harness and model show minimal looping and near-perfect accuracy
  (Table~\ref{tab:gpqa-loop-subdomain}).

\item \textbf{Cognitive PD default under uncertainty.}
  In the absence of tool grounding, the model defaults to principle-based
  deduction, the dominant cognitive mode for science questions.
  On tasks where the model is uncertain, PD rate rises from 0.52 (correct) to
  0.64 (regression), while IC (commitment) halves from 0.05 to 0.03 and RR
  (recovery) collapses from 0.018 to 0.005.
  The PD$\to$PD self-transition probability reaches 0.85: deduction begets
  deduction.

\item \textbf{Harness architecture without a commit deadline.}
  OpenClaw's open loop provides no structural signal to transition from
  deliberation to commitment.
  The model can re-enter PD indefinitely, and as the trajectory lengthens,
  entropy collapses from 0.24 (head) to 0.07 (tail) and 87\% of tokens fall
  into the low-entropy regime.
  The context budget is exhausted, and the answer extractor recovers nothing
  (89\% malformed, 83\% missing boxed answer on regressions).
\end{enumerate}

Removing any one condition prevents the spiral.
ZeroClaw removes condition (3) via its fixed template, absorbing the same PD
over-expression within a bounded budget.
OpenCode partially removes condition (1) by providing a tool-execution path that
forces mode transitions, and partially removes condition (3) by not permitting
unbounded re-entry.
DirectLLM removes all three by having no loop at all, achieving 80.3\% accuracy
with zero cap events.

The corrective implication follows directly: a harness that either provides an
executable tool path for every task class or imposes a structural commit
deadline after a bounded deliberation budget will eliminate context-exhaustion
failures without requiring any improvement in the base model's reasoning
capability.
